\documentclass[11pt]{article}
\usepackage{xcolor}
\usepackage{url}
\usepackage{amsmath, amsfonts, amssymb, amsthm}
\usepackage{mathtools}
\usepackage{verbatim}
\usepackage{graphicx}

\newcommand{\rem}[1]{\textcolor{red}{[#1]}}
\newcommand{\ed}[1]{\textcolor{blue}{#1}}
\newcommand{\ct}[1]{\rem{#1 --ct}}
\newcommand{\avar}{Fatom (AppUnintr 0 [])}
\newcommand{\bvar}{Fatom (AppUnintr 1 [])}
\newcommand{\cvar}{Fatom (AppUnintr 2 [])}
\newcommand{\fnot}[1]{Fnot (#1)}
\newcommand{\intr}[1]{\mathbb{I}_{\nu}(#1)}

\begin{document}
	\title{Invertibility Condition Proofs}
	\author{Arjun Viswanathan}
	\date{}
	\maketitle
	
\section*{Proof: Invertibility Condition Over BV Multiplication and Equality}
Equivalence to prove:
$$\forall s, t\ :\ \texttt{BV}_n.\ (-s\ |\ s)\ \&\ t = t
\iff (\exists x\ :\ \texttt{BV}_n. x \cdot s = t) 
$$
	
\subsection*{Case: $s = 0$}
If $s = 0$, we can prove both implications as follows:
\begin{itemize}
	\item Given $(-s\ |\ s)\ \&\ t = t$, we can derive
		\begin{align*}
			(-s\ |\ s)\ \&\ t = t &\models (-0\ |\ 0)\ \&\ t = t\\
			&\models 0\ \&\ t = t\\
			&\models 0 = t
		\end{align*}
		Since $t = 0$, we can prove the existence of $x$ by 
		picking $x = 0$ so that
		$$x \cdot s = 0 \cdot s = 0 = t$$
	\item Given $x$ such that $x \cdot s = t$, we have
		\begin{align*}
			t &= x \cdot s\\
			  &= x \cdot 0\\
			  &= 0
		\end{align*}
		Then, $(-s\ |\ s)\ \&\ t = (-0\ |\ 0)\ \&\ 0 = 0 = t$.
\end{itemize}

\subsection*{Case: $s \neq 0$}
At least one of $s$'s bits must be one. In other words,
$$s = s_{n-1}s_{n-2} \ldots s_{k+1}10 \ldots 0$$
There are $k$ least significant $0$s in $s$ where $k \geq 0$. Then,
\begin{align*}
	-s &= \neg s + 1\\
	&= \neg s_{n-1} \neg s_{n-2} \ldots \neg s_{k+1} 0 1 \ldots 1\\
	&= \neg s_{n-1} \neg s_{n-2} \ldots \neg s_{k+1} 1 0 \ldots 0
\end{align*}
and $(-s\ |\ s) =$\\
	\begin{tabular}{c c c c c c c c c c}
		& &$\neg s_{n-1}$ & $\neg s_{n-2}$ & $\ldots$ & $\neg s_{k+1}$ & $1$ & $0$ & $\ldots$ &$0$\\
		&$| $&$s_{n-1}$ &$s_{n-2}$ &$\ldots$ &$s_{k+1}$ &$1$ &$0$ &$\ldots$ &$0$\\
		&$= $&$1$ & $1$ & $\ldots$ & $1$ & $1$ & $0$ & $\ldots$ &$0$
	\end{tabular}\\
Then, $(-s\ |\ s)\ \&\ t =$\\
\begin{tabular}{c c c c c c c c c c}
	&$ $&$1$ & $1$ & $\ldots$ & $1$ & $1$ & $0$ & $\ldots$ &$0$\\
	&$\& $&$t_{n-1}$ &$t_{n-2}$ &$\ldots$ &$t_{k+1}$ &$t_k$ &$t_{k-1}$ &$\ldots$ &$t_0$\\
	&$= $&$t_{n-1}$ & $t_{n-2}$ & $\ldots$ & $t_{k+1}$ & $t_k$ & $0$ & $\ldots$ &$0$
\end{tabular}\\

Thus, we have:
\begin{equation*}
	\label{eq1}
	(-s\ |\ s)\ \&\ t = t_{n-1}\ t_{n}\ \ldots\ t_{k+1}\ t_k\ 0\ \ldots\ 0
\end{equation*}

Hence the invertible condition $(-s\ |\ s)\ \&\ t = t$ is equivalent to the equality \\
	\begin{tabular}{c c c c c c c c c c}
		&&$t_{n-1}$ & $t_{n-2}$ & $\ldots$ & $t_{k+1}$ & $t_k$ & $0$ & $\ldots$ &$0$\\
		&$= $&$t_{n-1}$ &$t_{n-2}$ &$\ldots$ &$t_{k+1}$ &$t_k$ &$t_{k-1}$ &$\ldots$ &$t_0$
	\end{tabular}\\
i.e. that the least significant $k$ bits of $t$ are $0$.

We now reduce to a statement about integers modulo $2^n$. Observe that we may write

\[
\text{bv2int}(s) = 2^{n-1}s_{n-1} + 2^{n-2}s_{n-2} + \cdots 2^{k+1}s_{k+1} + 2^k = 2^ks'
\]
for some odd integer $s'$. Similarly,

\[
\text{bv2int}(t) = 2^{n-1}t_{n-1} + 2^{n-2}t_{n-2} + \cdots 2^{k+1}t_{k+1} + 2^kt_k + 2^{k-1}t_{k-1} + \cdots + t_0
\]
is divisible by $2^k$ in $\mathbb{Z}$ (or equivalently in $\mathbb{Z}/2^n$) if and only the least significant $k$ bits of $t$ are $0$. We saw above that this is equivalent to the invertibility condition $(-s\ |\ s)\ \&\ t = t$.

On the other hand, $(\exists x\ :\ \texttt{BV}_n. x \cdot s = t)$ is equivalent to $(\exists x\ :\ \mathbb{Z}/2^n. x \cdot \text{bv2int}(s) = \text{bv2int}(t))$, since $\text{bv2int} : \texttt{BV}_n \cong \mathbb{Z}/2^n$ preserves multiplication.

Writing $s, t$ for $\text{bv2int}(s), \text{bv2int}(t)$, we are reduced to the following:

\emph{Let $s, t : \mathbb{Z}/2^n$. Suppose $s = 2^k s'$ where $s' : \mathbb{Z}/2^n$ is (the image of) an odd integer. Then $2^k$ divides $t$ in $\mathbb{Z}/2^n$ if and only if $s$ divides $t$ in $\mathbb{Z}/2^n$.}

This is basic number theory:

($\Leftarrow$) Suppose $s$ divides $t$. Let $x$ be such that $x s = t$. From the calculation

\[
 (xs')2^k = x(2^ks') = xs = t,
\]
we conclude that $2^k$ divides $t$.

($\Rightarrow$) Suppose $2^k$ divides $t$. Let $x'$ be such that $x' 2^k = t$. Since $s'$ is (the image of) an odd integer, it is relatively prime to $2^k$. By a result of number theory, $s'$ admits a multiplicative inverse $u$ in $\mathbb{Z}/2^n$. We compute

\[
(x'u)s = (x'u)(2^ks') = (x'2^k)(us') = x'2^k = t,
\]
so that $s$ divides $t$.

\end{document}
